\chapter{Is My Online Self Still Me?}
\section{A Profile Picture Kept for Three Years}
Pick up something: not your phone, but the smaller object on its screen, your profile picture.

It may be a cat, an anime character, a selfie cropped to half a face, scenery, or the default gray figure. Less than a centimeter across, it attends every occasion beside your name: group messages, friend requests, game results, late-night likes on classmates' posts. Teachers identify you through it; online friends who have never met you imagine you through it. Some people keep an avatar for three years, until friends can no longer separate the image from them. Seeing a similar cat outdoors, a friend photographs it and sends, ``This looks like you.''

Pause and feel how strange that sentence is: a cat looks like you.

The picture is only the beginning. You have a username, an online-only way of speaking, a carefully coordinated game skin, a fandom side account, perhaps a secret account for confessions, a tree hollow that keeps what you tell it. Together they form an online you, speaking, befriending, arguing, and receiving abuse on your behalf.

Does this you, assembled from avatar, username, and posting history, count as you? Wholly or partly? Must the eating, sleeping person offline answer for its mistakes? From this sub-centimeter square, we will travel to one of humanity's oldest questions: who am I?

\section{At Ten, Lime Starts Drawing}
Now enter a scene.

Time: Thursday at ten at night. Place: a curtained bedroom on the eleventh floor of an ordinary apartment building in a southern city. A desk lamp lights one corner of the desk; the milk tea has gone cold. Outside, an electric scooter honks once, then quiet returns.

A fourteen-year-old girl sits at the desk. At school, on registers, and in parents' WeChat groups, she is Chen Xiaotang: above-average grades, few words, and quiet written year after year in teachers' reports. Now, wearing headphones and holding a digital stylus, she draws a mechanical bird standing on ruins.

But nobody in this online art community knows Chen Xiaotang. They know Lime: an artist with thirty thousand followers, regular updates, answers to every question in the comments, and occasional late-night complaints. Her profile says, ``Just drawing. Thanks for liking it.'' Lime uses expressive conversational particles, leaves thoughtful comments on new followers' art, and advertises commissions by private message. Offline, Chen Xiaotang blushes even when raising her hand in class.

At 10:40, she finishes and posts the bird. Within ten minutes, likes climb: three hundred, eight hundred, fifteen hundred. Comments bustle: ``Divine composition!'' ``Master, need another pet cat?'' ``This made me cry. It reminded me of last year...'' Lime---Chen Xiaotang---replies to each, fingers moving with an ease unlike her daytime self.

At 11:30, Mom opens the door: ``Still not asleep?'' She answers and closes the app. The screen dims; Lime disappears. The room and girl remain unchanged. At seven tomorrow, Chen Xiaotang will shoulder her schoolbag and resume being the quiet, middling student.

Notice what happened tonight: not one person pretending to be another, but one person living genuinely for several hours in two worlds. Both laughs are real, and so are both kinds of tiredness.

\section{Which Is the Real Her?}
State the conflict without rushing to conclude.

Two opposite accounts are possible, each firmly believed by many adults.

First: the online version is fake. Lime is Chen Xiaotang's packaged image, showing only her best pictures and most agreeable words, without even a photograph of her face. The gap between offline introversion and online liveliness is acting. You know its extreme forms: distrust every online persona; behind a screen, who knows what is opposite you? Here the online self is a coat to remove at the door. Mistaking coat for body makes someone a gullible fool.

Second: the online version is the real one. Teachers' scrutiny, parental expectations, and classmates' judgments force Chen Xiaotang to fold herself away offline. Only Lime's account frees her from being the quiet, middling student. Birds, words, and friends are her choices. If someone feels truly themselves only before sleep, does their real self live by day or late at night?

Do not choose yet. A deeper question hides beneath online authenticity:

\textbf{What makes someone the same person?}

The same body? Its cells have been replaced countless times since you were seven, yet you remain you. The same name? Chen Xiaotang and Lime name one body, and people can have several names. Continuous memory? Then who inhabits a drunken blackout? Others' perceptions? School sees quiet; the internet sees a master. They see two versions.

This is no philosophical pastime; it has practical teeth. Can Chen Xiaotang refuse commissions Lime promised followers? If someone impersonates Lime to steal money, whom should police investigate? If an interviewer ten years later finds Lime's late-night complaints, is that normal background research or invasion of privacy? Before answering, clarify the relationship between the online and offline selves.

Choose a provisional position: is Lime Chen Xiaotang's false mask or true face?

\section[Removing Masks and Needing Them]{Those Who Remove Masks and Those Who Need Them}
The sharpest dispute over online identity occurs not in philosophy books but around a practical question: should anonymous internet use be allowed?

\textbf{Position one: show names and people will restrain themselves.}

Make this case fully. Its voices often come from profound wounds and deserve to be heard.

A cyberbullying victim's mother might say: thousands of accounts attacked my daughter with words too filthy to repeat. Their users then deleted the accounts, changed avatars, and carried on. Would they dare if every statement bore an identifiable person? There are at least three arguments. First, responsibility needs an address. Offline assault brings consequences because police can find the assailant; online harm flourishes without a return address. Second, anonymity makes courage unequal: hidden attackers dare everything while exposed victims do not even know their identities. Third, real names are an ancient social institution. Villagers hesitate to do wrong because everyone knows them and reputation follows for life. Real-name rules restore village constraints online.

This is a forceful argument. Anyone mobbed in comments understands the anger of losing a night's sleep while opponents do not even have names.

\textbf{Position two: masks protect people who speak truthfully.}

The other side has a full case, supported by history.

In late eighteenth-century America, thirteen newly independent states fiercely debated ratifying a new constitution. Hamilton, Madison, and Jay wrote eighty-five essays defending it under one shared pen name, Publius, a Roman hero's name. Collected as The Federalist Papers, they remain political classics. Why pseudonyms? They wanted readers to judge arguments rather than take sides according to authors' identities; pseudonyms also made criticism of those in power easier. A pen name did not falsify the essays, but clarified their reasoning.

Look closer to home. An isolated eighth grader, afraid to tell parents or teachers, writes late at night on a confidential account that they cannot go on and receives strangers' comfort and practical suggestions. An employee reporting workplace misconduct needs an anonymous letter. A child with an unshareable worry can ask whether they are normal only behind a pseudonym. The same mask is a weapon in some hands and an oxygen mask in others. Ask not whether masks are good, but who pays more if all are forbidden.

One country conducted a major legal experiment on this question, with unexpected results.

South Korea went further than most toward online real-name requirements. From 2007, large websites had to verify users' identities before allowing posts. The rationale followed the first position: a return address for every statement should restrain malice. Five years later, in 2012, the Constitutional Court ruled the system unconstitutional and abolished it. Its review revealed two things. Malicious comments had not markedly declined: determined abusers found workarounds while rule-following people fell silent first. Meanwhile, websites' collection of tens of millions of identities created a treasure trove for hackers, followed by major leaks. It was like packaging the nation's identity numbers and placing them where thieves most wanted to look.

This is an easily misjudged counterexample. Intuition equates anonymity with license and real names with civility. South Korea's experience is more tangled: malice disguises itself, truth retreats, and identity records become a new hazard. Both universal-real-name remedies and absolute-anonymity claims may stumble over this case.

Japan offers another picture. Its 2channel, founded in 1999, long ranked among the world's largest anonymous forums, covering anime through social criticism with almost no signed posts. It generated vicious comments, but also became a remarkably creative seedbed of Japanese internet culture, producing expressions, subcultures, and pointed social debate. Anonymity was fertilizer and poison together; cutting it out removed both.

No universally accepted answer exists. China's approach is a compromise. The Cybersecurity Law effective from 2017 requires platforms to verify identities for services including information publication, while users may display nicknames: real names backstage, voluntary naming onstage. A mask is allowed, provided an identifiable person stands behind it. Whether this compromise is good is among the questions you must weigh later.

\section[Frontstage, Backstage, and Fallen Walls]{Thinking Bridge: Frontstage, Backstage, and Fallen Walls}
Is there a portable tool for asking which self is real, without relying on intuition? Yes, but first consider an ordinary offline scene.

At a mall with classmates on Saturday, you use one voice, set of jokes, and way of walking. Answering parents' questions at home that evening, you use another tone and account. Nobody explicitly taught the switch; you have done it since learning to walk. Neither version is false. Relaxed with classmates and restrained with parents, you remain genuinely you before different audiences.

In the 1950s, sociologist Erving Goffman developed such observations into The Presentation of Self in Everyday Life. His central metaphor is theater: we perform \textbf{frontstage} before audiences in classrooms, dinners, offices---anywhere people watch and rules apply. \textbf{Backstage} lies beyond the audience's view: your bedroom or time with closest friends, where you remove makeup, complain, and rehearse. Often misread as declaring everyone a hypocritical actor, the theory instead treats performance not as deception but as social life itself. Being polite to a teacher is not falseness; the student role carries those lines. Without them, school could not function for a day.

The theory supplies three questions:

\begin{itemize}
\item Who is this performance's audience? Am I frontstage or backstage?
\item Which script and costume am I using now?
\item Could these different audiences meet?
\end{itemize}
The third is crucial. Traditional life separates audiences: your boss does not meet childhood friends; Grandma does not hear your trash talk on the court. Each circle watches its own play without entering another's performance. Call this \textbf{audience segregation}. That wall lets you be different selves without falling apart.

One of the internet's great acts was blowing up the wall.

Online, teachers, parents, classmates, strangers, and your own future self ten years hence may occupy one auditorium. Your teacher might discover a fandom side account; a future interviewer might unearth an extreme comment shared three years ago. This is \textbf{context collapse}: previously separate family, school, friendship, and stranger contexts fall into one room. Unsure who watches, you no longer know which script to use. The public-private boundary is not merely deliberately crossed; technology melts it.

Apply the three questions to Chen Xiaotang. Her performances work because the art audience never meets school and family audiences. If a classmate discovers Lime's identity, the wall falls. The problem ceases to be which self is true and becomes practical sociology: where should she stand when two scripts play simultaneously before one audience?

Use the tool tomorrow. Before posting, ask whom you imagine watching, who might actually watch, and whether your statement holds up if those audiences meet. This is not slickness but guarding the door for your future self.

\section[Names: Ancient Testimony]{Name as Shadow or Baggage: Ancient Testimony}
Now read a brief primary source. Ancient thinking about names and actual people was no shallower than ours.

In Zhuangzi's ``Free and Easy Wandering,'' Yao offers rule of the world to the recluse Xu You, who refuses. Among his reasons is:

\begin{quote}
Name is the guest of reality.
\end{quote}
Broadly, name---reputation, standing, title---is subordinate to reality, the actual person or thing, as a guest depends on a host. The host comes first; the guest follows. Xu You means the title of ruler may sound honorable but is not what he actually wants. Why pursue an accessory?

Beside our question, this is strikingly sharp. Over two thousand years ago, Xu You identified a system of names able to detach from reality, reproduce, drift, and supplant its host. Usernames, avatars, follower counts, and personas are names; the eating, sleeping person who hurts and tires is reality. Zhuangzi warns against seating the guest in the host's place: do not let Lime's thirty thousand followers decide what Chen Xiaotang draws, says, or becomes.

But historical evidence points the opposite way too. Chinese people may be among the world's most practiced at living under several names. Traditional scholars had personal names, courtesy names, and literary sobriquets, sometimes several changing with circumstance. Su Shi's courtesy name was Zizhan and his sobriquet Dongpo Jushi, Layman of the Eastern Slope, from land he farmed during exile in Huangzhou. Demotion, relocation, or a changed outlook might bring another name. In modern times, pen names became writers' standard equipment: Lu Xun was Zhou Shuren's most-used pseudonym among hundreds over his lifetime. Nobody considered Su Shi fraudulent or Lu Xun split into personalities. Different names marked different facets: the broad-minded Eastern Slope layman and the exacting memorial-writing Su Shi together formed the full reality.

Together, these halves suggest something more interesting than whether online selves are true or false: \textbf{the internet did not invent separation between name and reality; it accelerated and mass-produced it.} The question is old; the scale is new. An ancient writer prepared ink and paper for a new sobriquet; you change an alias in ten seconds. Old names circulated in handwriting; your sentence can be copied ten thousand times in an afternoon.

\section{One Abusive Comment, Three Explanations}
Now compare explanations. Begin with an undisputed account of what happened.

An eleventh-grade boy is mild offline, shy with strangers, and silent when someone cuts ahead in the cafeteria. In a team-based competitive game, however, teammates know him as a toxic ranter. One mistake triggers a screenful of sarcasm and cutting words, with shouted voice chat audible down the hallway. Asked whether he behaves that way offline, he pauses: ``That's just online. It isn't really me.''

Separate four things. One person uses different words and actions in two settings: fact. He says the online version is not really him: his own account. Whether the abuse is justified: a value judgment, set aside briefly. We compare the middle layer: \textbf{why does this happen?} At least three explanations differ, each with reasons and gaps.

\textbf{Explanation one: anonymity removes restraints, or disinhibition.}

Psychologists describe an online disinhibition effect. Without showing faces or names, with others invisible and intangible, restraints loosen: shame, awareness of expressions, anticipation of consequences. Some drinkers become talkative not because alcohol contains another person, but because it loosens the brakes. Here the boy has not become bad; absent offline constraints---faces, onlookers, costs of offense---allow existing cruelty to leak out.

This concise explanation fits much experience and helps explain nastier comment sections. Its gap is equally clear: it explains released cruelty, not released kindness. Why do some anonymous users become harsher while others become gentler and more willing to share, like the child seeking confidential help? If anonymity merely releases brakes, why do some cars head toward kindness? The explanation makes online environments too uniform.

\textbf{Explanation two: the platform directs the performance. Conditions shape expression.}

Change viewpoint. Ask not first what is wrong with a person, but which behavior an environment rewards. Each platform has an invisible directing mechanism. Character limits simplify complexity, and simplicity can become harshness. Likes and shares reward emotional intensity, sinking mild truths and lifting heated partial truths. Followers and rankings turn visibility into competition, making expression increasingly performative. Recommendations gather similar viewpoints, and even moderate people become extreme amid allies' applause. The boy's cruelty is therefore not solely exposed nature but learned behavior. Abusing teammates provides immediate rewards: shifting blame, asserting status, venting before an audience. Gentleness receives none.

This sharp account locates new responsibility: ask designers about their buttons rather than merely scold users. Its limit is making people puppets when everything belongs to platforms. Some players in the same game do not abuse anyone. Environments shape but do not make the final decision. Assigning all responsibility to algorithms is as lazy as assigning it all to human nature.

\textbf{Explanation three: he was never only one self.}

This goes further: perhaps asking which version is real is mistaken. Recall Goffman. People have different frontstages for different audiences; mildness faces offline observers, anger faces teammates, and both are authentically him. His denial may reverse the truth: the online version is also him, formerly backstage without an audience. Personas are not false skins pasted over real people, but existing facets granted separate stages. \textbf{Virtual identity is not false identity}; it is a real personality appearing in a particular theater.

This best accommodates what the others miss: online good and evil can both be real because the self is multiple. But beware its easiest misuse: excuse-making. Calling abuse merely an online persona can sound airtight here, yet the person reduced to tears suffers real harm. Several facets do not mean none is responsible. Explanation traces behavior's origins; it does not issue an automatic pardon. Confusing explanation with responsibility is the evasion this chapter warns against.

Together, disinhibition explains loosened restraint, platform theory explains guiding rules, and multiple selves explain human branching. These are different magnifications, not mutually exclusive options. Distrust claims that one lens suffices; they usually discarded facts beforehand.

\section[Theseus's Ship and Your Lost Password]{Deeper Challenge: Theseus's Ship and Your Lost Password}
Now for an ancient intellectual tool seemingly prepared for today's question.

In the Life of Theseus within Parallel Lives, Greek biographer Plutarch records a puzzle. Athenians preserved Theseus's ship, replacing decayed planks one by one until none was original. Philosophers disputed whether it remained Theseus's ship.

The question survives over two thousand years because it concerns identity: after every component changes, why is something still the same? Apply it to people. Cells turn over, opinions repeatedly reverse across a decade, experiences alter personality. Why are you at fourteen the same person as at four? Further, Lime and Chen Xiaotang share a body but scarcely share memories---followers know nothing of her daytime life---relationships, or names. Are they the same person?

Seventeenth-century English philosopher John Locke offered an influential answer in An Essay Concerning Human Understanding: personal sameness lies not in the body but in \textbf{continuity of consciousness}, plainly put, memory. Remembering yesterday's deeds and reasons ties you backward through a chain. Where the chain breaks, you break. A person with a drunken blackout needs others to explain last night's actions: not merely embarrassment, but a small philosophical accident, with behavior ownerless in the sense of memory.

Locke's useful test throws sparks when shone on the internet age.

First, memory is uneven. Locke assumes memory resides in minds, fading and disappearing, with forgetting a human default. The internet outsources it to servers. You forget posts from three years ago; servers remember. You think deletion ends them; screenshots and archives remember. A monster Locke never imagined appears: \textbf{your past remembers you better than you do}. Theseus's puzzle reverses. The person whose planks keep changing is pinned by imperishable old planks---past statements---and made accountable for every one.

Second, chains break. Forget an account password and fail to recover it, and posts, friends, and memories become an unclaimed estate. Are those still the person's words? They vaguely remember saying them but not why, to whom, or with what feelings. The chain breaks not in time but in \textbf{context}: the words survive, their world does not. This helps explain the injustice felt when old posts are excavated for public judgment. A statement is ripped from its original ship and installed on today's for trial.

Besides body and memory, is there a third rope binding personal sameness? Yes: \textbf{recognition}, perhaps especially relevant online. Others keep identifying you: family use a name, friends recognize a face, society follows one set of records. The self is largely a collective project, half your narration and half others'. The internet separates the project: school says quiet, followers say master, family says well-behaved. Three files, three selves, each with witnesses. Theseus's question becomes: when three groups claim authority over the real you, who decides?

Philosophers have no standard answer. But take a steadier approach: instead of asking which self is real, ask to whom each version can answer. Identity may be less a true-or-false question than a craft requiring continual renewal. Each day, you reclaim the selves scattered across different rooms.

\section[Internet Memory and Loneliness]{An Unforgetting Internet and Unconnected Loneliness}
This ancient problem now bites through three developments.

\textbf{First: forgetting has become a right.}

For most human history, forgetting was free and automatic. Move cities and childhood embarrassments disappear; time removes witnesses to youthful mistakes. The internet canceled that default. In a landmark 2014 case, a Spanish citizen asked the European court to require Google to remove links to newspaper notices over a decade old, concerning a property auction connected with debts long since settled. The court supported his central claim: under certain conditions, search engines should remove links to outdated information. The judgment spread the phrase right to be forgotten: may someone prevent their past permanently heading search results? The EU's GDPR, effective in 2018, subsequently specified an erasure right; China's Personal Information Protection Law, effective in 2021, also allows deletion requests in particular circumstances.

Law answers only half. The everyday half involves audience groups, three-day-visible posts, deleting updates one by one before closing an account, and hoping old classmates never searched embarrassing histories. These features share a desire not to lie but for a \textbf{right to change}: to prevent the past self permanently mortgaging the present. Xu You's guest-and-host formulation has a modern version: reality keeps growing while names freeze on servers.

\textbf{Second: emotions travel through network cables.}

In 2014, an academic paper disclosed a controversial experiment. Researchers including Facebook data scientists adjusted positive and negative content proportions in roughly 690,000 users' feeds for one week without their knowledge, then observed their posts. Users shown less positive material posted slightly gloomier messages; those shown less negative material posted slightly more brightly. Effects were small but consistent. Emotions could cross screens without handshakes, meetings, or even awareness of another person's existence.

The dispute concerned procedure, not the conclusion: those roughly 690,000 people never knew they were subjects. Supporters said adjusting feeds was ordinary daily recommendation work. Opponents replied that routine practice does not justify experimenting on people, especially to manipulate emotion. The dispute remains, but fixes one fact in place: your emotions are not entirely yours; some come from your feed.

\textbf{Third: unprecedented connection and unprecedented loneliness arrived together.}

Online communities are real, whatever adult prejudice says. Rare-disease family groups can tell you at three in the morning where to find medicine. Far-flung gaming guildmates remember your birthday. Comments beneath Lime's art, recalling someone's painful previous year, are real people recognizing one another. Declaring all online friends fake is as lazy as declaring every persona false.

Yet other signals point the opposite way. In 2018, Britain appointed a minister for loneliness, treating isolation as a public problem requiring national action; Japan created a similar position several years later. Unprecedented connectivity and officially recognized loneliness need not contradict each other. As connection becomes cheap, its quality stratifies. As every emotion finds an amplifier, gathering for warmth and gathering for anger use identical equipment. You have probably seen both.

Together, these developments restate the old question. When names never vanish, emotions spread remotely, and communities are always present, how can the reality at the center---the eating, sleeping, hurting you---retain the host's seat?

\section{For You: A Delete Button}
Return to Chen Xiaotang's bedroom.

Imagine a realized technology: on your eighteenth birthday, every account presents a button. Pressing it permanently deletes everything you posted as a minor across the internet, including servers, caches, and search engines, irretrievably in every form. Decline, and everything remains in your public record. Law gives the choice solely to you, not your parents.

The button now faces eighteen-year-old Chen Xiaotang. She remembers Lime's thirty thousand followers, late-night encouragement, and first commission payment. She also remembers incoherent posts during a breakdown at fourteen, screenshots exposing others during arguments, and writing that now makes her blush.

Press or not?

There is no standard answer. Before deciding, place everything from this chapter on the scales.

Name and reality can separate: deleting the name may protect the growing reality. But does it? Memory binds personal sameness: are forgotten posts still hers? Cyberbullying victims ask whether one-click deletion also pardons those who harmed others. Confidential help-seekers ask whether they would ever have spoken without a future delete button. Platforms direct expression through likes and algorithms: how much embarrassing history did she write, and how much did the environment write through her hands?

Is deleting the past renewal or escape? Is retaining everything honesty or punishment? What way back should a fair society leave a fourteen-year-old poster?

Answer with reasons. Notice that every reason for whether she should delete may someday become another person's measure for deciding whether to forgive your past self.
